\documentclass[12pt]{article}
\usepackage[utf8]{inputenc}
\usepackage[margin=0.75in]{geometry}
\usepackage{amsmath}\usepackage{amssymb}\usepackage{enumitem}\usepackage{multicol}\usepackage{tcolorbox}
\setlength{\parindent}{0pt}
\begin{document}

\begin{center}{\Large \textbf{Identify maximums, minimums, and where a function increases or decreases}}\\[2pt]{\small Pre-Cal \textbar\ CK-12: Maximums and Minimums}\end{center}\vspace{0.25cm}

{\large\textbf{Key Idea}}\par\smallskip
A \textbf{maximum} is a highest point; a \textbf{minimum} is a lowest point. A function \textbf{increases} where the graph rises left-to-right and \textbf{decreases} where it falls.\par


{\large\textbf{Key Terms}}\par\smallskip
\begin{itemize}[leftmargin=1.4em]
  \item \textbf{Extremum} — a max or min point
  \item \textbf{Interval} — the stretch of $x$-values where a behavior holds
\end{itemize}


{\large\textbf{Worked steps}}\par\smallskip
Step 1: Find the turning point (vertex or peak/valley).\par
Step 2: Read increasing/decreasing on each side.\par
Example: $y=x^2$ has a \textbf{minimum} at $(0,0)$; it decreases for $x<0$ and increases for $x>0$.\par


{\large\textbf{You Try}}\par\smallskip
Does $y=-x^2$ have a max or a min, and where?\par

\end{document}